\clearpage

\section{Conditional E-variables for 2x2 tables}
\label{sec:conditional-2x2}



\subsection{Definitions and calculation methods}

\subsubsection{Conditional E-variable}

The FNCH distributional setup is collected in Section~\ref{sec:fnch}.
Conditioning on the total number of ones,
and using the pdf from Eq.~\ref{eq:fnch_pdf},
\begin{align}
 	S(x \mid \delta,n_a,n_b,n_1) 
 	&= \frac{f(x\mid \delta,n_a,n_b,n_1)}{f(x\mid 0,n_a,n_b,n_1)} \\
 	&= \exp\{\delta x + A(0; n_a,n_b,n_1) - A(\delta; n_a,n_b,n_1)\}
\end{align}


As we do not know $\delta$, we marginalized wrt some prior on $\delta$,
which is the most computational intensive part.
$$
	S(x \mid n_a,n_b,n_1) = \int_{\Rf} S(x \mid \delta,n_a,n_b,n_1) p(\delta) \dd{\delta}
$$

\subsubsection{Two Bernoulli random variables}

\emph{Heavily adapted on Dante's Thesis.}

Let $X,Y$ be two independent Bernoulli random variables,
\begin{equation*}
	X \sim \mathrm{Bernoulli}(\theta_a) \quad
	Y \sim \mathrm{Bernoulli}(\theta_b).
\end{equation*}

$\Pc$ describes the set of distributions on $(X,Y)$ 
where $X,Y$ are independent and marginally distributed as Bernoulli.
A distribution on $(X,Y)$ could be parameterized by
$P_{\theta_a,\theta_b}$ where $\theta_a,\theta_b$ ranges in $[0,1]$.
There's a \emph{natural} way to re-parameterize the tuple via log-odds ratio,
\begin{equation*}
	\delta:= \log \frac{\theta_a}{1-\theta_b} \Big/ \frac{\theta_b}{1-\theta_b} \in \Rf.
\end{equation*}
Note that two parameters (a tuple) are still needed,
the log-odds ratio $\delta$ is chosen such that 
we can inexplicitly represent a composite set, i.e.
\begin{equation*}
	P_{\delta} = \{P_{\theta_a,\theta_b} \mid \log(\frac{\theta_a}{1-\theta_b} \cdot \frac{1-\theta_b}{\theta_b}) = \delta \in \Rf \}.
\end{equation*}

Assuming point null and alternative hypothesis,
it is easy to formulate the testing problem by a likelihood ratio test,
\begin{align*}
	\Hc_0: X,Y \sim P_{\theta_a,\theta_b} \quad
	\Hc_1: X,Y \sim P_{\theta_a^*,\theta_b^*}.
\end{align*}

The problem extends easily to sequential test in iid settings.
\begin{align*}
	\Hc_0: (X_1,Y_1),(X_2,Y_2),\dots \iid P_{\theta_a,\theta_b} \quad
	\Hc_1: (X_1,Y_1),(X_2,Y_2),\dots \iid P_{\theta_a^*,\theta_b^*}.
\end{align*}

% TODO: the convexity of the distributions on X, Y and X, Y

\subsubsection{Setup and notation}

We denote the Kullback–Leibler divergence by $\KL{P}{Q} = \int \log(\dv*{P}{Q}) \dd{P}$
where $\dv*{P}{Q}$ is the Radon-Nikodym derivative of $P$ with respect to $Q$ and
the integral is taken on the same probability space (not shown here).
We will assume that the measures discussed here all have full support.

We start with a binary random variable $U$
following Bernoulli distribution with parameter $\mu$,
here $\mu$ is also the mean in mean-value parameterization.
The density reads
\begin{align*}
	p_{\mu}(U)
	 & := \mu^U \cdot (1 - \mu)^{1-U}, \quad U \in \{0,1\}, \quad \mu \in [0,1]
\end{align*}

By the standard properties of regular exponential families,
we parameterize $\Pc$ in terms of the mean-value parameter space $\Mt_p$ and write
\begin{align*}
	\Pc = \{P_\mu: \mu \in \Mt_p\}, \quad \Mt_p = [0,1].
\end{align*}

The canonical parameterization for any $\mu \in \Mt_p$ is defined as
\begin{align*} 
	p_{\beta}(U) = \frac{1}{Z_p(\beta)} \cdot \exp(\beta \cdot U) \cdot p_{\mu}(U), 
\end{align*}\label{eq:bern_exp}
where $Z_p(\beta)$ is the normalizing constance and the canonical parameter space is
defined as $\Bt_p = \{\beta: Z_p(\beta) < \infty \} = \Rf$.
% \beta = \log(\frac{\mu}{1-\mu})

Moving on to the 2-stream examples with different group size $n_a, n_b$.
At each time point $i$, $n_a$ and $n_b$ copies of Bernoulli variables are collected,
following the same distribution:
\begin{align*}
	U_a & = (U_{a,1}, U_{a,2},\cdots,U_{a,n_a})^T \sim \mathrm{Ber}(\bm{\mu}_a) \\
	U_b & = (U_{b,1}, U_{b,2},\cdots,U_{b,n_b})^T \sim \mathrm{Ber}(\bm{\mu}_b)
\end{align*}
The sufficient statistics vector is denoted by $X = (\sum_j^{n_a}U_{a,j}, \sum_j^{n_b}U_{b,j})$.
$\Pc$ is extended by the i.i.d. assumption to
a set of distributions for random process $U_{(1)}, U_{(2)}, \dots$
for $i \geq 1$, i.i.d. copies of $U_{(i)}=(U_{(i),a}, U_{(i),b})$.
Using superscript, we denote the random process up to time $i$ by
$U^{(n)} := (U_{(1)}, U_{(2)}, \dots U_{(n)})$
and $X^{(n)} := (X_{(1)}, X_{(2)}, \dots X_{(n)})$ for the sufficient statistics respectively.

The \emph{alternative} we consider first $\Qc = \{Q\}$ is simple with some oracle
$\bm{\mu}^* = (\bm{\mu}_a^*, \bm{\mu}_b^*)$.
\begin{align*}
	\hypA: U_a \sim \mathrm{Ber}(\bm{\mu}_a^*), U_b \sim \mathrm{Ber}(\bm{\mu}_b^*) \quad
	\text{independent}.
\end{align*}

Under mean-value parameterization:
\begin{align*}
	\hypA: X_a \sim P_{\bm{\mu}_a^*}, X_b \sim P_{\bm{\mu}_b^*} \quad
	\text{independent}.
\end{align*}

The null distribution $\Pc$ is the set of distributions where group $a$ and $b$
are equipped with the same parameter and hence composite:
\begin{align*}
	\hypN: X_a \sim P, X_b \sim P \quad
	\text{i.i.d. for some } P.
\end{align*}


\subsection{Generated exponential family}

This is based on 4.3 of the simple case and the example 1 and section 4.1 of the general case paper.

Pls refer to the \eqref{eq:bern_exp}, here the null is again z

\subsubsection{Convexity prerequisite}

\begin{remark}

The Bernoulli model, $\Pc=\{P_\theta:\theta\in[0,1]\}$, on $\Yc=\{0,1\}$ is convex, 
whereas the product model $\Pc^n=\{P_\theta^n:\theta\in[0,1]\}$ on $\Yc^n$ is not convex for $n>1$.

\end{remark}

\begin{proof}
	Convexity is straightforward for single observation case $Y=\{0,1\}$.
	For $\lambda\in [0,1]$, consider the convex combination of two elements,
	$P_{\theta_1},P_{\theta_2}\in \Pc$.
	Probability for $\lambda P_{\theta_1}+(1-\lambda)P_{\theta_2}$,
	\begin{align*}
		\PR{Y=1} &= \lambda \theta_1+(1-\lambda)\theta_2 \\
		\PR{Y=0} &= 1-\bigl(\lambda \theta_1+(1-\lambda)\theta_2\bigr)
	\end{align*}
	Therefore, the mixture is again in $\Pc$ with parameter $\theta^\circ=\lambda \theta_1+(1-\lambda)\theta_2\in[0,1]$.
	
	Take $n=2$, consider the mixture
	\begin{equation}
		P = \frac{1}{2} P_0^2 + \frac{1}{2} P_1^2
	\end{equation}
	where $P_0^2$ put all its mass on $(0,0)$ and $P_1^2$ put all its mass on $(1,1)$.
	
	However, there is no $\theta\in[0,1]$ for $P_\theta^2$ such that,
	\begin{align*}
		&P_\theta^2(0,0) = P(0,0) = \frac{1}{2} 
		&P_\theta^2(1,1) = P(1,1) = \frac{1}{2} \\
		&P_\theta^2(1,0) = P(1,0) = 0
		&P_\theta^2(0,1) = P(0,1) = 0
	\end{align*}
\end{proof}


\textbf{Why finite convexity does not immediately imply closure under arbitrary mixtures.}

Let
$$
\mathcal P=\{P_\theta:\theta\in\Theta\}.
$$

Convexity of $\mathcal P$ means that for every finite collection
$$
P_{\theta_1},\dots,P_{\theta_k}\in\mathcal P
$$

and weights
$$
\lambda_1,\dots,\lambda_k\ge 0,
\qquad
\sum_{i=1}^k \lambda_i=1,
$$

we have
$$
\sum_{i=1}^k \lambda_i P_{\theta_i}\in\mathcal P.
$$

However, an arbitrary mixture has the form
$$
P_W=\int P_\theta\,dW(\theta),
$$

where $W$ may have infinite support.

Convexity alone directly handles finite sums such as
$$
\sum_{i=1}^k \lambda_i P_{\theta_i}.
$$

It does not directly imply that
$$
\int P_\theta\,dW(\theta)\in\{P_\theta:\theta\in\Theta\}.
$$

A simple example showing the issue is the open interval
$$
\mathcal C=(0,1)\subset\mathbb R.
$$

The set $\mathcal C$ is convex. Every finite convex combination of points in $\mathcal C$ remains in $\mathcal C$. However, convexity alone does not imply that limits of such convex combinations must remain in $\mathcal C$.

For example, take
$$
x_m=\frac{1}{m}.
$$

Then $x_m\in\mathcal C$ for every $m$, but
$$
\lim_{m\to\infty} x_m=0.
$$

The limit point satisfies
$$
0\notin\mathcal C.
$$

Thus convexity alone does not imply closedness under limits or arbitrary infinite averaging. Therefore, to conclude that an arbitrary mixture belongs to the model, one needs an additional argument.

In the finite sample-space case, this additional argument is supplied by Carathéodory's theorem.

\textbf{Why the integral is in the convex hull.}

Assume the sample space $Y$ is finite:
$$
Y=\{y_1,\dots,y_m\}, m \in \Nf.
$$

Then each distribution $P_\theta$ can be represented by the vector
$$
p_\theta
=
\bigl(p_\theta(y_1),\dots,p_\theta(y_m)\bigr)
\in\mathbb R^m.
$$

The mixture distribution $P_W$ has probability mass function
$$
p_W(y)
=
\int p_\theta(y)\,\dd{W(\theta)}.
$$

Therefore its vector representation is
$$
p_W
=
\left(p_W(y_1),\dots,p_W(y_m)\right).
$$

Substituting the definition of $p_W(y_i)$ gives

\begin{align}
p_W
&=
\left(
\int p_\theta(y_1)\,dW(\theta),
\dots,
\int p_\theta(y_m)\,dW(\theta)
\right) \\
&=
\int
\bigl(p_\theta(y_1),\dots,p_\theta(y_m)\bigr)
\,dW(\theta) \\
&=
\int p_\theta\,dW(\theta).
\end{align}

This is an average, or barycenter, of the points
$$
\{p_\theta:\theta\in\Theta\}\subset\mathbb R^m.
$$

Therefore
$$
p_W\in\operatorname{conv}\{p_\theta:\theta\in\Theta\}.
$$

Since this convex hull is finite-dimensional, Carathéodory's theorem implies that there exist
$$
\theta_1,\dots,\theta_k\in\Theta
$$

and weights
$$
\lambda_1,\dots,\lambda_k\ge 0,
\qquad
\sum_{i=1}^k\lambda_i=1,
$$

with $k\le m+1$, such that

$$
p_W
=
\sum_{i=1}^k \lambda_i p_{\theta_i}.
$$

Thus the arbitrary mixture $p_W$ can be represented as a finite mixture.

If the model $\{P_\theta:\theta\in\Theta\}$ is convex, then this finite mixture is again equal to some single model distribution:

$$
\sum_{i=1}^k \lambda_i P_{\theta_i}
=
P_{\theta^\circ}
$$

for some $\theta^\circ\in\Theta$.

Therefore

$$
P_W=P_{\theta^\circ}.
$$

\subsubsection{Bayesian mixture update}

Assume the data generating process means are fixed, group sizes are also fixed,
we constructed a e-process if we \emph{hack} and know the oracle log-odds ratio,
at time $t$,
\begin{equation*}
	S(B_t) = S_{\textsc{cond}, \delta}(B_t) = \frac{p_{\textsc{fn}}}{p_{\textsc{hyp}}}
\end{equation*}
where $p_{\textsc{fn}}$ is the Fisher's noncentral hypergeometric mass
and $p_{\textsc{hyp}}$ is the central hypergeometric mass.

Cumulative product is worse linearly compared to GRO when sample size increase,
\begin{equation*}
	S^n_{\textsc{gro}} - \prod_1^{n} S_{\textsc{cond}, \delta}(B_t) \approx O(n) 
\end{equation*}
while the `oneshot' COND is approximately GRO
\begin{equation*}
	S^n_{\textsc{gro}} - S_{\textsc{cond}, \delta}(B^n) \approx O(1) 
\end{equation*}

Since we do not know $\delta$, we try a bayesian mixture
\begin{equation*}
	S(B_t) = S_{\textsc{cond}, W}(B_t) = \int S_{\textsc{cond}, \delta} \dd{W}.
\end{equation*}

Let's say we do not update $W$, always use the same $W$,
cumulative evidence is
\begin{equation*}
	\prod_{i=1}^t \int S_{\textsc{cond}, \delta} \dd{W}.
\end{equation*}

\textbf{Can this be simplified?}

But that's boring, we need to update the prior density,
\begin{equation*}
	\int S_{\textsc{cond}, \delta} \dd{W_{t-1}}.
\end{equation*}

How do I update this

\subsubsection{Bayesian prior calculations}

\paragraph{Integral over prior}

\paragraph{Gaussian prior}

Suppose we choose a standard Gaussian prior on $\Rf$
$$
    \delta \sim \Nc(\mu_0,\sigma_0^2).
$$

After observing one table, with conditional observation $Y_a=y_a$, the posterior density is
$$
    p(\delta\mid y_a,n_a,n_b,n_1)
    \propto
    \exp\left\{
        \delta y_a
        - A(\delta)
        - \frac{(\delta-\mu_0)^2}{2\sigma_0^2}
    \right\}.
$$

This posterior is generally not Gaussian and does not have a standard closed form.

For multiple tables,
suppose we have observed $B_i=(y_{a,i},n_{a,i},n_{b,i},n_{1,i})$, $i=1,\dots,t$
$$
	p(\delta\mid B^t) \propto 
	\exp\left\{
		\delta\sum^t_{i=1} y_{a,i} 
		- \sum^t_{i=1} A_i(\delta)
		- \frac{(\delta - \mu_o)^2}{2\sigma_0^2}
	\right\}
$$

\paragraph{Conjugate prior}

In the literature, conjugate prior for general exponential family can be written:
\begin{align*}
	p(\delta) \propto \exp\{n_0x_0 \delta - n_0 A(\delta)\}.
\end{align*}
where $x_0$ is suggested to use the expectation of $x$ under $\delta=0$,
which is hypergeometric distribution so $x_0=\frac{y_a\cdot y_b}{n}$.

After seeing one data point, assume $n_a,n_b,n_1$ fixed,
the posterior is
$$
    p(\delta\mid y_a,n_0,x_0)
    \propto
    \exp\left\{
        (n_0x_0+y_a)\delta - (n_0+1)A(\delta)
    \right\}.
$$

It should be noted that the conjugacy only works for same margin,
else the cumulant function would change.

\todo{then what is the point of conjugate then?}


The \emph{conditional \E-value} with a fixed $\delta$ is
\begin{equation*}
	S_{n_a,n_b,n_1,\delta}(y_a) = \frac{f_{\textsc{FN}(n_a,n_b,n_1,\delta)}(y_a)}{f_{\textsc{FN}(n_a,n_b,n_1,0)}(y_a)}.
\end{equation*}


The `Bayes' is w.r.t. a (truncated) Gaussian with zero mean and scale $\delta$,
\begin{equation*}
	\int_{\Rf} S_{n_a,n_b,n_1,\delta}(y_a) \cdot \pi(\delta) \dd{\delta}
\end{equation*}

See: \url{https://github.com/wyq977/e-table/blob/yongqi/src/e_table/methods/bayes.py}

\subsubsection{IID and convexity}

In the original testing framework,
both null and the alternative require sample to be distributed iid,
i.e.,
For every \emph{binary random variable} $X$ in group a or group b,
at time $j\in\Nf$,
$$
X_{a,1}^{(j)},X_{a,2}^{(j)},\dots,X_{a,n_a}^{(j)}\iid \mathrm{Bernoulli}(\theta_a).
$$

We observe two streams of binary random variables
$$
X_{a,i}^{(j)}, X_{b,i}^{(j)} \in \{0,1\},
$$
where
\begin{itemize}
	\item $g \in \{a,b\}$ denotes the group,
	\item $j \in \Nf$ denotes the batch/time index,
	\item i denotes the within-batch observation index.
\end{itemize}


It is also true that at the coarse filtrations (batch level),
$X$ also following the same Bernoulli.

That is to say, for stream $a$, 
we are always tossing the same coin over and over again in time,
and we just observe them in the coarse filtration of $n_a$ coins.

The finest filtration is alway two infinite streams of Bernoulli,
i.e., two coin-flipping experiments.

\begin{enumerate}
	\item Now, null is tossing the same coin in group a and b,
	alternative is tossing two different coin in group a and b.
	But the coins remains the same throughout.
	We collect them at a coarse filtration, fixed, $n=n_a+n_b$.
	\item Still the same null and alternative,
	could I change the collections (coarsening)?,
	i.e., every time it's not guaranteed $n_a,n_b$ in group a and b.
	I could have stopped and observed $5,10$ in the first time,
	and observed $100, 20$ in the second time.
	\item Assuming we have the same coarsening,
	i.e., we observed same $n_a,n_b$ throughout.
	Within each observation block, we use the same coin (iid),
	but between blocks, we changed coins.
	\item I am not sure how to formulate this one,
	maybe the null is fixed two coins and the alternative is anything but,
	coarsening variable.
\end{enumerate}

\newpage

\subsubsection{E-zoo comparison}

\paragraph{UI}

Based on the general case paper and the UI slide.

Let $\hat{\mu}_{\mid i-1}$ be any \emph{non-anticipating}\todo{What does it mean?}
estimator based on the first $i-1$ observations and we
denote the MLE estimator for the null as
$\hat{\mu}_{\mid n} = \argmax_{\mu\in\Theta_0} p_{\mu}(U^{{(n)}})$.

\begin{align*}
	S_{\textsc{UI}} := \frac{\prod_{i=1}^n q_{\hat{\mu}_{\mid i-1}} (U_{(i)})}{p_{\hat{\mu}_{\mid n}}(U^{(n)})}
\end{align*}

By i.i.d. assumption, the denominator reads the likelihood of $U^{(n)}$ given the MLE estimator.

It is mentioned that \emph{prequential ML} method is used here.

% TODO: verify whether the MLE estimator can also be used in the numerator.

\paragraph{RIPr}

By Theorem 1 in \cite{turnerExactAnytimevalidConfidence2023}, the GRO \E-variable
reads:
\begin{align*}
	S_{\textsc{RIPr}} := \frac{p_{\mu^*}(U^{(n)})}{p_{\mu^\circ}(U^{(n)})},
\end{align*}
where the prior put mass 1 on the boundary and $\mu^\circ = (n_a\mu^*_a/n + n_b \mu^*_b/n)$.

\begin{remark}
	In 2-stream Bernoulli and the null hypothesis is 2 group is distributed with the same distribution $P \in \Pc$,
	$S_{\textsc{cond}} = S_{\textsc{RIPr}} = S_{\textsc{PSEUDO}} = S_{\textsc{MIX}}$,
	which achieves GRO by corollary 1 of safe testing paper.
\end{remark}

\paragraph{COND}

% TODO: add a section on Wald's ratio test.

We can first write the conditional \E-variable for a data block $X = U = (U_a, U_b)$:

\begin{align*}
	S_{\textsc{COND}} := \frac{q(X_a, X_b \mid X_a + X_b)}{p(X_a, X_b \mid X_a + X_b)}
	= \frac{q(X_a\mid X_a +X_b)}{p(X_a\mid X_a+X_b)}
	= \frac{f_{\textrm{FN}(n,n_a,Z,w)}(X_a)}{f_{\textrm{hyp}(n,n_a,Z)}(X_a)}.
\end{align*}

\begin{remark}
	Right now, I finally realize 4.2.3 of Yunda's thesis where the \E is:
	\begin{align*}
		S_{\textsc{COND}} := \frac{p_{\bm{\mu}}(X^{k-1}\mid Z)}{p_{\spri{\mu}}(X^{k-1}\mid Z)}.
	\end{align*}
	Here the $Z:=\sum_{j=1}^k X_j$ is the sum of the sufficient statistics and
	$X^{k-1}$ would be the sufficient statistics losing $1$ \emph{degree of freedom}.
\end{remark}

In our special case of Bernoulli streams, the conditional distribution corresponds to
the Fisher's noncentral hypergeometric (FN) distribution and the denominator is
simply the cases when odds ratio $\omega = 1$ (hypergeometric distribution).

\begin{align*}
	S_{\textsc{\tiny COND}} = \prod_{i=1}^n \frac{q(X_a^{(i)}\mid Z^{(i)})}{p(X_a^{(i)}\mid Z^{(i)})}
	= \prod_{i=1}^n \frac{f_{\textrm{FN}(n,n_a,Z^{(i)},w)}(X_a^{(i)})}{f_{\textrm{hyp}(n,n_a,Z^{(i)})}(X_a^{(i)})}
\end{align*}

\clearpage

\subsection{Questions I already answered}

\subsubsection{Prior or turning distribution}

We assume the oracle is some fixed point $\mu^*$
Let $\pi$ be the prior distribution on the alternative distributions.
The prior predictive distribution is defined as
\begin{align*}
	p(U) = \int p_{\mu}(U) \pi(\mu) \dd{\mu}.
\end{align*}

In the discussion with Ly on 02.05.2025, he pointed many names for
the above expression:
maximal likelihood distribution, the prior predictive distribution (Bayesian), and Bayes mixture (MDL).

Additionally, he mentioned that we usually used \textsc{letters} ($X,A$) for
observables where \textsc{Greek letters} are reserved for unknown
($\pi, \theta, \Theta$), especially in Bayesian framework.

In contrast to the belief's $\mu$ in Bayesian, MDL principle
regards $\pi$ on $\mu$ rather as a turning tool, uninformative,
to approach the oracle point $\mu^*$ given incoming data.

\subsubsection{What is effect size and why do we need it?}
\cite{turnerGenericEvariablesExact2024} deals with effect size by stating that $\hypA$ consists of hypotheses with a minimal effect size $\delta$ while the parameter space of interest is $\hypN(\delta)$ (null hypothesis spaces) in \cite{turnerExactAnytimevalidConfidence2023}. These has confused me a great deal but I will try to explain in my words:

The effect size is the minimal difference between group a and group b
we would like to see when rejecting the null hypothesis.
Since we are concerned with a parametric model of Bernoulli stream in both group,
parameterized by $\theta_a$ and $\theta_a$.
The effect size would essentially be the divergence $d$ between $\theta_a$ and $\theta_b$
in the parameter space.

\subsubsection{Why put a prior on the boundary only}

$W_1^*$: Prior on $\Theta(\delta)$; $W_1'$: Prior over $\Theta_1^>(\delta)$.
What are good reasons of putting all prior mass on the boundary?

Pros:
\begin{itemize}
	\item Computationally simpler as we only calculate once on $\theta_a$
	      and we get $\theta_b = \theta_a + \delta$ for free.
	\item If $(\theta_a^*, \theta_b^*)$ sits somewhere near the line,
	      \E-variable growth rate is higher with larger prior support.
	\item Theorem 2 in \cite{turnerExactAnytimevalidConfidence2023} guarantees a degenerate prior
	      on the boundary $\Theta(\delta)$
\end{itemize}

Cons:
\begin{itemize}
	\item If $d_{abs}(\theta_a^*, \theta_b^*) \gg \delta $,
	      \E-value calculated from $W_1'$ is much larger than $W_1^*$
\end{itemize}

\sout{What does it mean to \emph{pick $\hypA$ restricted to $\Theta(\delta)$}?}

In the discussion with Peter, he mentioned that the project $P^*$ might not be on the diagonal line unlike section \ref{sec:case1}.
It might just be a highly non-convex function that we can only solve by finding the min. of KL divergence after discretization.
He mentioned something about the projection not being straight but rather curve.

I think I made a mistake here, when looking at the divergent case,
the null would be the divergence bigger than $\delta$ while the alternative would be else?

{\color{red}
	can we plug in $\theta_b^\circ = \theta_a^\circ + \delta$ here? If so,
	it would just be solving a cubic equation of $\theta_a^\circ$.
}

\subsubsection{Extended null and one-shot conditional e-variable}

The extended null hypothesis is meaningful now!

In this setup that I call ``one-shot'' where the whole data $U_a, U_b$ is seen at once,
the conditional \E-variable is optimal.

The extension of the null hypothesis is meaningful, in this parametric Bernoulli family only:
the extended null $\hypN'$ will be all distributions where all permutation $U$ are the same
and the marginal distributions of all (all exchangeable sequences).

In other words, let $U'=(U_1,U_2,\dots,U_{n_a+n_b})$ and
$\pi$ be any permutation operator on $\{1,\dots,n_a+n_b\}$.
The null hypothesis is now extended to a much bigger one,
namely, whether two streams $U_a, U_b$ are exchangeable,
\begin{align*}
	\hypN':=(U_1,U_2,\dots,U_{n_a+n_b}) \overset{d}{=}
	U_{\pi(1)},U_{\pi(2)},\dots,U_{\pi(n_a+n_b)}\}.
\end{align*}

Moreover, the marginal distribution of $U_j \overset{d}{=} U_i$ for any $i, j$
but the conditional dependence is not assumed in $\hypN'$ any more
in contrast to independence assumed in the original $\hypN$
where $U_i, U_j \iid P$.

Nevertheless, the $S_{\textsc{cond}}$ is still GRO-optimal, we can claim it is a
\textit{robust} \E-variable for a bigger problem.

\subsubsection{Sequential-ness and conditional-ness}

I was confused about the formulation between the sequential ratio and
the (what I called) global conditional \E-variables.

Let $X,Y$ be Bernoulli variables with $p$ and $q$, then we can write the
conditional likelihood ratio test (Wald's) as follow:

\begin{align*}
	R^n := \prod_{i=1}^n \frac{f_1(X_i, Y_i \mid X_i + Y_i)}{f_0(X_i, Y_i \mid X_i + Y_i)}
\end{align*}

and $f_1(\cdot) := (1 + \theta)^{-1}$ and $f_0 = (1 + \theta^{-1})^{-1}$

Here is the example of composite testing:
\begin{align*}
	\hypN: (X_1, Y_1), (X_2,Y_2),\dots \iid P \in \Pc_1 \: \textrm{v.s.}
	\hypA: (X_1, Y_1), (X_2,Y_2),\dots \iid P \in \Pc_{\theta}
\end{align*}

where $\Pc_1$ and $\Pc_{\theta}$ are distribution parametrized by odds ratio $\theta$.

Now it is almost the same as before but just extended where each data point $(X_1, Y_1)$
is replaced by a contingency table.
I think the reason I got so confused is here:
we are testing $m$ copies of tables $T_i$ and within that $T_i$,
each table is a 2x2 table of $n_a, n_b$ Bernoulli streams with $p, q$.

Using the notation $X^b_i, i = 1,\dots, n_a$ Ber and $X^a_i, i =1,\dots,n_b$
The full expression reads:
\begin{align*}
	\frac{f_1(X_1^a, X_2^a,\dots X_{n_a}^a, X_1^b, X_2^b,\dots X_{n_b}^b \mid Y_a + Y_b)}
	{f_0(X_1^a, X_2^a,\dots X_{n_a}^a, X_1^b, X_2^b,\dots X_{n_b}^b \mid Y_a + Y_b)}
\end{align*}

For a particular table $T_i$, I can just jump to the common doc, the hypothesis being tested is

\begin{quote}
	For all the 2x2 tables $T_i$, group a and group b frequency
	has no difference.
\end{quote}

\begin{align*}
	\hypN: T_1, T_2, \dots \iid P \in \Pc_1
	\quad \textrm{v.s.} \quad
	\hypA: T_1, T_2, \dots \iid P \in \Pc_{\theta}
\end{align*}

Questions to discuss tomorrow:

\begin{itemize}
	\item Is the \E-process-ness still here? \emph{YES}, by
	\item Is the optional continuation still here, I think so
	\item For the mixture, $\theta$ gets Bayesian updating each time (numerical)
	\item I think we assume group size $n_a, n_b$ is stationary right? Does the
	\item I remember there is some nice properties concerning the concave of the Fisher's information
\end{itemize}

\subsection{Questions without answers}

\subsubsection{Conditioning, margins, and coarsening}

% FIXME: this subsection partially overlaps with "Setup and notation" above.
% Consider merging once the notation and problem statement are stabilized.

\begin{quote}
	This is based on the talk with Wouter after e-reader.
\end{quote}

Start simple, assume two streams of binary random variable $X,Y$ indexed by $t$.
If they are both iid Bernoulli, then we are in the paired example where each time we threw two coins.
At each time, $X_t$ or $Y_t$ is observed, which is governed by some process or switch.
For simplicity we also assume this process to be Bernoulli iid with parameter $\eta>0$, i.e.
$p(I_t=1) = \eta$ (seeing a action from group a).
It assumed independent for simplicity now and note that the means are stationary and iid.
\begin{align*}
	P(X_t \mid I_t) \iid \mathrm{Bernoulli}(\theta_a) \\
	P(Y_t \mid I_t) \iid \mathrm{Bernoulli}(\theta_b)
\end{align*}

\begin{table}[h]
	\centering
	\begin{tabular}{lrr}
		\toprule
		    & $a$ & $b$ \\
		\midrule
		$1$ & 0   & 1   \\
		$0$ & 0   & 0   \\
		\bottomrule
	\end{tabular}
\end{table}

Let the process goes on for $n$ rounds, and we collected the data until time $n$ into the summarized table.
\begin{table}[h]
	\centering
	\begin{tabular}{lrrr}
		\toprule
		    & $a$      & $b$      &       \\
		\midrule
		$1$ & $n_{a1}$ & $n_{b1}$ & $n_1$ \\
		$0$ & $n_{a0}$ & $n_{b0}$ &       \\
		    & $n_a$    & $n_b$    & $n$   \\
		\bottomrule
	\end{tabular}
\end{table}

In plain english, we are interested in the presence of the differences between two streams.
$\eta=1$ corresponds to the paired-sample settings where we \emph{must} observe two binary variables (two coins).
In other words there's one more degree of freedom by allowing $\eta\in(0,1)$.

\begin{itemize}
	\item $\Hc_0: X_1,X_2,\dots,X_n\iid \mathrm{Bernoulli}(\theta_0)$ and $Y_1,Y_2,\dots,Y_n\iid \mathrm{Bernoulli}(\theta_0)$
	\item $\Hc_1: X_1,X_2,\dots,X_n\iid \mathrm{Bernoulli}(\theta_a)$ and $Y_1,Y_2,\dots,Y_n\iid \mathrm{Bernoulli}(\theta_b)$
\end{itemize}

Let $\delta$ be some measure of divergence between $\theta_a$ and $\theta_b$.
Conditioning on $n_a,n_b,n_1$ gives the FNCH distribution
from Eq.~\ref{eq:fnch_pdf}.

Instead, we can just condition on the total number of ones
% TODO: complete this expression or turn it into prose.
\begin{align*}
	P(n_{a1} \mid n_a+n_b,n_1,\delta) =
\end{align*}

The questions remains are:
\begin{itemize}
	\item Without knowledge on specific $n_a,n_b$,
	      we can still formulate this into a likelihood ratio test between $\delta$ and $\delta=0$,
	      but I think this is composite null vs composite alternative.
	      This is easy to see, let $r=\frac{n_a}{n_b}\in(0,1)$,
	      the likelihood ratio with the above formulas must contain $r$.
	      And $r$ could take value in $(0,1)$, making the null and alternative composite.
	\item Is it natural to also condition on $n_a$ on top of $n=n_a+n_b, n_1$?
	      % FIXME: incomplete note from discussion with Wouter.
\end{itemize}


A quick word on the Group Invariance paper, we are testing Gaussian with any variance.

\begin{align*}
	 & \Hc_0: X \sim \Nc(0, \sigma)  \quad \sigma \in \Rf         \\
	 & \Hc_1: X \sim \Nc(\mu, \sigma), \quad \mu, \sigma \in \Rf.
\end{align*}
The variance is a nuisance parameter that testing simple does not care.
In other words, we can normalized or scale the data however we like.
The question is: the conditioning on $n_1$ is in the same spirit
and we implicitly were also doing the same conditioning $n_a$ (and $n$)
so that the composite null \emph{collapse} into a single distribution.

\subsubsection{Is MIX also an e-process?}

Adopting the same notation as in Yunda's thesis,

\subsubsection{Why Turner's e-variable for the first table is 1}

% TODO: answer this question.

\subsubsection{Finding confidence sequence for case 1}
First of all, can we find CS for this testing setup.
If so, the confidence set should be searched through the $\Theta_0$
(dashed red line in \refFig{fig:case1_GRO}) or the entire $[0,1]^2$?

\subsubsection{Convexity of hypothesis spaces}

I am not sure why the modified settings where \cite{turnerExactAnytimevalidConfidence2023} and \cite{turnerGenericEvariablesExact2024} looks at the single outcome first and then extend to a single data block.

Calculating the KL divergence over the set of distributions $\mathcal{H}$ are convex while the calculation of KL over parameter space $\Theta$ is \textbf{ALSO} convex??

\begin{quote}
	% \textit{
	note that we cannot draw this conclusion if $\hypN'$ is not convex; for then the distribution $p_{W^\circ}'$ may not correspond to the distribution $p_{W^\circ}$ in the original problem — this correspondence is only guaranteed if $p_{W^\circ}'$ coincides with some $p_{\theta}'$.
	% }
\end{quote}

The part I am most confused about is the relationship between the convexity of the hypothesis space (original and modified) and parameter space. What we want is a $\hypN'$ that is compact and convex.

\href{https://statproofbook.github.io/P/kl-conv.html}{Proof: Convexity of the Kullback-Leibler divergence}, it gives the proof of convexity in the pair of probability distributions $(p, q)$ in discrete cases but the same could be applied to continuous cases with integration.

\subsubsection{General case paper discussion with Peter}

If we have 1-d regular exponential family with the same sufficient statistics,
let's assume we are talking about Turner's example.

The GRO one is the conditional \E-variable and putting the prior on the curve (ark)
results also in the RIPr one (GRO) \E-variable.

Compared to this case (fixed odds ratio $\delta$ vs $1$, curve vs diagonal straight line),
it is not clear that if we still get the optimal \E-variable if instead of a fixed $\delta$,
we put a prior on the simplex, what sort of optimality we would get if there any.

Would the conditional variable still be the optimal one?


Let $f(k) = \binom{n_a}{k} \binom{n_b}{U-k}$ and $f'(k) = \frac{n_b!}{k!(n_a-k)!(U-k)!(n_b-U+k-n_a)!} = \binom{n_a - U + k}{n_a}$


\begin{align*}
	p_{n_a, n_b, u, \omega}(k) = \frac{\frac{n_b!}{k!(n_a-k)!(u-k)!(n_b-U+k-n_a)!}\omega^k}{\sum \frac{n_b!}{k!(n_a-k)!(u-k)!(n_b-U+k-n_a)!}\omega^j}
\end{align*}


We know $n_{11}$ is the $Y_a = k$ in our settings,

\subsection{Other misc}

\subsubsection{Project direction on 28.07.2025}

The goal to get a conditional \E-variable
\begin{itemize}
	\item with priors on log odds ratio: Jeffrey's, conjugate, Cauchy, Laplace, Gaussian
	\item uniformly most powerful (UMP) test for a fixed $\alpha$, ask Sebastian
	\item normalized maximal likelihood (NML), aligns with the MDL principle
\end{itemize}

The second amazingly achieves the most powerful test with a given test,
related materials are in the reply to the safe testing but it is given in
the context of unconditional one. So Peter was not sure about its look here.

Compare it with the old one (unconditional one)
\begin{itemize}
	\item Turner's
	\item Universal inference (MLE)
	\item p-to-e calibrator ($e = 1/\sqrt{q} - 1$ maybe)
\end{itemize}

Previous results shows that the UI performs the worst, next to Turner's,
a simple calibrator works best, e.g. $f(p) = \kappa p^{\kappa-1}$ for some $\kappa \in(0,1)$
and $f(p) = p^{-1/2}-1$.
noted that the Fisher's exact test gives a very small p-value

First step is to start with a fixed prior and fixed alternative.

\subsubsection{Re-orientate at 01.04.2025}

The biggest reason that I have been so disoriented is that I have so many documents, tests
scattered across platforms and places: \githubIssue, \personalRepo, \forkedRepo.
Originally, I wanna use Issue as a tracking board, personal repo for experiments,
and the forked repo for dev.
In the end, I became over-whammed by all and often found myself wasting energy,
organising thoughts even after just 1 day.

Thus, I have decided that I will put everything related to 2x2 in this \LaTeX\ document.
I will still use my \personalRepo for testing but I will not use \githubIssue any more.
I will write exclusively here as I find writing with \LaTeX\ is still the simplest
when even a little bit of math is involved.
I will use the item or numbered list here and \sout{strike through it} upon completion
or delete it when it is not longer relevant.
